\chapter[More Than Machines]{The Industrial Revolution Changed More Than Machines}
\section{A Nineteen-Yuan T-Shirt}
Pick up something you may be wearing: a pure-cotton T-shirt.

Online, it costs 19.90 yuan with free delivery. Small print on its label identifies where it was sewn: perhaps Vietnam, Bangladesh or an industrial district in southern China. It omits the full journey. Cotton may grow in Texas or Xinjiang, be mechanically compressed into enormous bales and shipped to another country's spinning mill. Yarn travels to a third country's knitting and dyeing factory, then a fourth country's workshop for cutting, stitching and branding before flying to your door.

A 19.90-yuan garment travels farther than you may in a lifetime, yet remains cheap enough to discard after one season.

Try a thought experiment: hand it to someone 250 years ago. In London, Beijing or Delhi in 1770, cotton clothing was property an ordinary household took seriously. Europe's poor often owned only one or two respectable outer garments, included in wills and passed down. It was not that people lacked thrift, but that cloth cost so much: between a cotton boll and a length of fabric stood thousands upon thousands of hand-working hours. A skilled Indian weaver might spend years producing the finest muslin. An English farmwoman turning her wheel between agricultural tasks might spin all day without enough yarn for gloves.

A century later, in 1870, British machines swallowed millions of cotton bales annually, producing enough cloth to measure in miles. People really began using miles for output. Cotton fabric became commonplace worldwide, so ordinary you can own a dozen T-shirts without noticing.

What happened in that century? Textbooks answer the Industrial Revolution, followed by a list: spinning jenny, Watt engine, railway, steamship. The list is correct but misleading: it makes machines seem the protagonists.

Machines were merely the most conspicuous actors. Timetables, childhood, households, urban rivers and skies, distant fields and mines whose inhabitants had never seen an engine were rewritten. Industrialisation rearranged a human day, a lifetime and humanity's relationship with the planet.

\section{Manchester, One Morning in 1842}
Come into a scene.

It is a winter morning in 1842, a quarter to five, in Manchester, north-western England: the world's most industrialised city, nicknamed Cottonopolis.

Naturally it is dark; even eight brings little winter daylight. Yet the city is awake, summoned by factory bells rather than sunrise. Riverside mill boilers have burned all night. A huge engine breathes on the ground floor, transmitting power through shafts and belts to hundreds of machines across five storeys.

A twelve-year-old girl, call her Annie, hurries barefoot over thinly iced paving. Shoes are saved for Sunday. Hundreds accompany her through unlit streets towards the mills. Gates close precisely at five; latecomers lose half a day's wages. Traffic or oversleeping is no excuse. Factory time is decided by clocks, not negotiation.

Enter the workshop and sound arrives before sight. Hundreds of machines make conversation possible only by shouting into ears. Older workers learn lip-reading, a skill unnecessary outside. Next come heat and damp. Dry air breaks cotton yarn, so steam keeps the room humid and windows remain closed. Fine fibres drift through warm, sticky air like golden mist in gaslight. Beautiful but lethal: years of work bring diseased lungs and lifelong coughs.

Annie is a piecer, immediately joining broken yarn on spinning machines. Her fingers must outrun the spindles. Smaller children crawl under running machinery to sweep accumulated fluff, the workshop's most dangerous task. Their bodies fit gaps adults cannot enter; machinery will not pause even half a second for an arm inside.

A floor below, stokers shovel coal into furnaces. It comes from mines dozens of kilometres away, where children far younger than Annie spend whole days underground without sun. Later you will hear one child's words.

At seven that evening, Annie finishes: fourteen hours. She returns to a back-to-back terrace in the workers' district, two households sharing a rear wall, without piped water. The street shares a pump and several toilets; sewage runs down an open central gutter to the river. Fish lived there thirty years earlier. Now coloured foam floats on it as upstream dye works empty leftovers: each day's orders determine the river's colour.

Remember this morning. Every later argument returns to Annie's feet, the workshop mist, boiler-house coal and colour-changing river.

\section{Machines Were Not New. What Was?}
Remove a misleading impression: industrialisation did not invent machines.

Over a millennium earlier, Song Chinese ironworks drove blast-furnace bellows with waterwheels. Ancient Romans and medieval monasteries milled grain with water. First-century Greek engineer Hero of Alexandria even made a steam-powered rotating sphere, a sound engine in principle, treated merely as a temple curiosity. Seventeenth-century Bengali weavers made muslin like flowing water, whose quality European machinery needed nearly a century barely to match.

The question was never why humanity suddenly learned to build machines. It always could.

Two other things were new.

First, machines began \textbf{reproducing themselves}. In 1764, Lancashire weaver Hargreaves made the spinning jenny, letting one person operate eight spindles, later over a hundred. Arkwright patented the water frame in 1769, producing coarse but strong warp yarn. Crompton combined them in the mule in 1779, making fine, tough yarn. Weaving immediately lagged behind the yarn supply, pressing for improved looms. Power looms appeared around 1785, then yarn again became insufficient and spinning improved. Each advance created another link's bottleneck, demanding a solution. Spinning and weaving raced for half a century like runners refusing to stop. Earlier machines were isolated inventions; these formed a \textbf{self-accelerating system}.

Second, machines escaped rivers. Water-powered mills needed fast-flowing streams, placing early factories in remote valleys. In the 1770s and 1780s, Scottish engineer Watt improved steam engines. Improved is the word: Newcomen's engines had pumped mines for decades, but were bulky, coal-hungry and confined to reciprocating motion. Watt's crucial contributions reduced fuel consumption and supplied steady rotary power. Topography no longer commanded energy: with coal, factories could locate anywhere. They converged on mines, ports and one another, reorganising cities around power.

Railways followed. In 1825, the first publicly passenger-carrying steam railway opened between Stockton and Darlington. Liverpool--Manchester followed in 1830, first proving that trains could carry freight and passengers profitably. Over twenty years, Britain laid more than ten thousand kilometres of track, distributing coal, iron, machines, newspapers, workers and punctuality. Railway companies demanded uniform national timetables; previously towns followed local church clocks, commonly differing by a quarter-hour.

These novelties reveal two difficult questions:

\textbf{First: why late-eighteenth-century Britain?} China had coal; machines existed widely; India and China had abundant skilled craftspeople; Dutch capital and markets were established. Why did Manchester become Cottonopolis?

\textbf{Second: what was this transformation?} One account calls it humanity's greatest liberation, making cotton clothing, rail travel and endless manufactured goods affordable. Another calls it a mincing machine: Annie's fourteen hours, child miners, dyed rivers, distant slaves and ruined weavers. Both hold real evidence. Can both be true?

\section{The People beside the Machines}
Manchester's morning appears differently through different eyes. Hear several voices before judging.

\textbf{Workers.} For Annie and her parents, factories first meant losing their time. Farm work was hard but seasonal, alternating frantic labour and respite according to weather. Factories followed clocks: over three hundred days yearly, long daily hours, fines for lateness, talking or opening windows. Rules covered workshop walls. Around 1810, English stocking-makers and cloth finishers reached their limit, gathering at night to smash new machines. Authorities called them Luddites, supposedly after the worker Ned Ludd. Troops entered industrial districts; machine-breaking temporarily became a capital offence.

Insert \textbf{a frequently misjudged counterexample}. Luddites acquired a reputation as ignorant mobs opposing progress. That is badly wrong. Most were experienced skilled workers, not enemies of machinery itself. They targeted particular machines in particular factories \textbf{using technology to cut wages, replace skilled adults with cheap children and pass inferior products off as craft goods}, sparing employers following accepted rules. They opposed distribution, not technology: owners received benefits while workers paid costs. The movement failed; its question outlived every machine it broke.

\textbf{Entrepreneurs.} Now fully present the factory owner's case. This is not forgiveness. Without understanding an opponent's reasons, no genuine argument can be won.

A cotton manufacturer in the 1820s might say: my investment came from relatives and friends, with my entire fortune at stake. Cloth prices change yearly, Liverpool cotton prices daily. If the neighbouring mill installs a larger engine, failure to follow means death. I did not seize children from schools: universal schooling does not exist. Their parents work here and the family needs every wage. Enclosure makes displaced country people still more desperate; they walk miles to my gates themselves. My factory pays Manchester taxes, and I funded a Sunday school. Machines neither strike nor drink, but coal costs shilling after shilling. How much profit do you imagine? The last cotton-price collapse closed three mills larger than mine.

Some of this is true. Early industrial Britain lacked social security; poor people already lived near subsistence. Employers risked fortunes and often failed. Factories cheapened goods, eventually benefiting consumers including the poor. But it evades a distinction. The investment risk was chosen; Annie's work was not. An illiterate twelve-year-old whose family needs her pay cannot meaningfully contract freely. \textbf{Not all forms of consent carry equal weight.} That is the defence's weakest point.

\textbf{Families.} A quieter change rarely appears beside machine lists: households were reorganised. Agricultural families produced together, children working beside parents and learning skills under a shared roof. Factories extracted work from home. Individuals earned separate wages; father and Annie entered different workshops or mines, sharing a home but spending days in different worlds. Women entered factories too, sometimes outnumbering men in cotton mills. Scandalised commentators deplored declining morals, seldom asking why women earned roughly half men's wages. Only later, in the mid-to-late nineteenth century, did restrictions on women's and children's work and rising men's wages spread the male-breadwinner, female-homemaker nuclear family among British workers. Notice the order: traditional family ideals did not simply shape industrial society; industrial society reinvented the traditional family. Today's familiar arrangement is itself partly an industrial product.

\textbf{The far end of the world.} Cottonopolis grew no cotton. In the early nineteenth century, Manchester chiefly consumed the American South's slave-grown crop. Sugar, rubber and coffee followed broadly similar stories. India also supplied cotton, despite its earlier leadership in textiles. Seventeenth-century Indian muslins sold globally, among the East India Company's most profitable goods. By the mid-nineteenth century, British machine-made cloth flooded India, bankrupting handweavers and converting a textile exporter into a raw-cotton exporter and Britain's largest cloth market. Governor-General Dalhousie's railway programme in the 1850s carried inland cotton seawards and British cloth inland. One track, two directions, two destinies.

Farmers across America, Africa and Asia were drawn in too. They had not seen steam engines, but steam engines saw them. Tightening world markets priced their cotton, indigo, tea and rubber; a Manchester fluctuation could starve villages half a globe away. Industrialisation happened across a map, not merely inside factories.

\section[Thinking Bridge: Three Conditions]{Thinking Bridge: Divide ``Why'' into Three Kinds of Condition}
Single answers to Britain's industrialisation---coal, steam or an innovative national character---are satisfying and usually wrong. Take a portable tool: \textbf{separate three kinds of condition}.

\textbf{Necessary: without it, the event cannot occur.} Ask whether removing it leaves the outcome possible. Without affordable, accessible coal, steam power was uneconomic and nineteenth-century industrial expansion unimaginable. Consider similarly cotton, suitable for machinery with near-limitless demand, or networks transporting food, materials and products.

\textbf{Sufficient: its presence is enough.} \textbf{No single factor sufficed for industrialisation.} Many places had coal. Eleventh-century Northern Song China offers a counterexample: large coal-fired iron industries near Kaifeng, with output some economic historians rank far above the contemporary world and impressive even against later Europe. Yet no Industrial Revolution followed. Coal, iron and advanced skills were insufficient. This wedges apart the habit of assuming X automatically generates Y.

\textbf{Triggers and catalysts: not foundations, but determinants of when, where and how fast.} Examples include high wages, discussed shortly; legal protection for patents and property; engineers and instrument-makers exchanging experimental results; or an invention clearing a particular bottleneck.

Why learn this? It treats single-cause worship. On hearing that someone succeeded through effort, a company failed through poor management or a relationship ended through incompatible personalities, ask three questions. Is effort necessary: do people succeed without it? Sufficient: does everyone trying succeed? Or catalytic: does it accelerate an outcome once other conditions exist? Most confident claims will hesitate. Hesitation is how thought begins.

Take this cause-separating blade to a contemporary source.

\section{An Eight-Year-Old's Testimony}
In 1842, a parliamentary commission reported on child miners. Britain had heard something of factory children, but nobody had really seen underground conditions. Testimony and illustrations shocked the public and directly prompted legislation banning women and children under ten from underground work that year.

One interview concerned eight-year-old Sarah Gooder in Yorkshire. As a trapper, she watched ventilation doors, opening them at the sound of coal wagons and closing them afterwards. Her testimony is paraphrased here:

\begin{quote}
I mind a trap in the pit and must stay there. I have no light; it is completely dark. In the daytime I am a little frightened because I cannot see. With a light I sing; in darkness I dare not sing---I dare not.
\end{quote}
(Source: Britain's 1842 Children's Employment Commission report on mines, witness Sarah Gooder, aged eight. The passage paraphrases its meaning.)

Notice its nature: not fiction or political essay, but a child's answers recorded word for word, carrying a diary's rough edges. Fear, singing and not daring are unadorned truths. State its limits too: commissioners selected and recorded testimony with reform agendas, and a child's memory and expression can distort. It establishes that such conditions and fears existed, not that every mine and child were identical. Recognise evidence's power and boundary together.

Place it beside another wall. The 1833 Factory Act broadly prohibited employing children under nine in cotton mills, limited nine-to-thirteen-year-olds to eight hours daily and first appointed salaried state factory inspectors. Eight hours seems enormous now; then it was unprecedented. The state first declared parental and employer agreement insufficient: children's working hours were law's concern. Previously freedom of contract meant voluntary agreement without state interference. After 1833, freedom acquired a lower bound.

Here two forces collided. One argued parents voluntarily supplied children for real wages, and prohibition would first starve poor households. The other said a society where an eight-year-old dared not sing in darkness had miscalculated. Real costs existed outside every ledger. Neither side won completely. Minimum standards advanced centimetre by centimetre: nine years, eight hours, inspectors, then ten years, ten hours and compulsory education. Each step provoked predictions of industrial collapse. Industry adapted and continued growing. Remember that pattern in today's arguments.

\section{Why Britain? Four Accounts, Four Blind Spots}
Historians have debated Britain's priority for over half a century. Four representative accounts are not ice-cream flavours, but different views of how history works, each with evidence and limits.

\textbf{First: wages and coal arithmetic.} Robert Allen argues that eighteenth-century British wages were among Europe's highest while coal was exceptionally cheap. Expensive labour plus cheap energy made \textbf{labour-saving invention profitable in Britain but not elsewhere}. Jennies saved labour and engines burned coal, as if designed for British prices. Lower French wages reduced incentives; Jiangnan's still lower wages and abundant workers reduced them further. This numerical account turns supposed national inventiveness into arithmetic. It explains rapid adoption and improvement better than initial sparks: the jenny's maker probably was not calculating Europe-wide wage tables.

\textbf{Second: institutions protected those taking risks.} Douglass North's school emphasises relatively early parliamentary checks, property protection and patent rules, with the 1624 patent statute a landmark. Inventors risked years expecting officials and powerful people not to seize results; merchants financed mills and railways because contracts were enforceable. Innovation required rules and backing, not mere daring. But compared with whom were institutions good? Eighteenth-century Britain also had corruption, electoral bribery and biased justice. Institutions themselves emerged historically: if they explain innovation, what explains them?

\textbf{Third: underground coal and overseas ghost acres.} Kenneth Pomeranz's Great Divergence argues that advanced regions such as England and Jiangnan remained similar into the eighteenth century, facing land and resource limits. Divergence came later. Britain had two advantages of fortune: coal near industrial districts and possession of New World resources. North American and Caribbean colonies supplied invisible ghost acres: slave-produced cotton, sugar and timber saved enormous domestic land and labour. Britain's own land could not have fed industrialisation. This rejects a chosen-Europe narrative in favour of extraction and contingency, while admitting coal and colonies explain \textbf{conditions}, not inevitability.

\textbf{Fourth: circulating knowledge.} Joel Mokyr stresses an Industrial Enlightenment joining understanding to practical skill. Scientists experimented, artisans built instruments, cafés and societies circulated knowledge, and strangers sought expert advice by letter. Steam improvements depended on atmospheric pressure and vacuum knowledge, not inspiration alone. This captures a soft but real ecology in which knowledge \textbf{reproduces}. Its limitation is measurement. Cultural climates cannot be weighed like coal; unmeasurable explanations readily become self-congratulation, with innovative culture explaining innovation circularly.

These accounts need not exclude one another, but are \textbf{not equivalent}. The first offers the most checkable figures; the third restores neglected victims; the second and fourth identify real but hard-to-measure factors. A deep historical skill is withholding premature allegiance while seeing \textbf{where each shines light and what it leaves dark}.

\section[Further Challenge: Paying for Cheapness]{Further Challenge: Who Pays for ``Cheap''?}
An economic tool explains industrialisation's deepest contradiction and much in your basket: \textbf{externality}.

Consider a transaction. You pay 19.90 yuan for a T-shirt, receiving clothing while the seller receives money. Both agree; the books are clean. Economists would say its price contains cotton, spinning, weaving, dye, sewing and transport costs accepted by both parties.

Return to Manchester's colour-changing river. Dye works paid nobody for dumping; downstream residents received no compensation for foul water and dead fish. Boiler smoke entered lungs without payment for coughing. No ledger recorded Sarah Gooder's dark childhood. Costs occurred but \textbf{fell on people outside the transaction}. These are negative externalities: prices withheld truth, and apparently cheap goods passed bills to absent people.

Externalities can be positive. Railways collect fares while nearby land rises in value, news spreads and opportunities appear, benefits they cannot charge for. Uncaptured benefits weaken incentives to build enough, prompting government subsidies for railways, roads and education. Literacy benefits everyone; parents alone cannot bear its bill.

The concept reframes growth coexisting with costs. Cheap cloth partly reflected not efficiency but \textbf{transferred bills}: downstream residents, workers' lungs, mining children, southern American slaves, ruined Indian weavers and ultimately the atmosphere. Coal's carbon dioxide knocks on nobody's door, accumulating quietly and billing people two centuries later. Nineteenth-century Britain could not know this. That does not excuse wrongdoing; it identifies externality's most insidious form, invisible even to those transferring the cost.

Use the tool carefully at two boundaries.

First, its framework assumes everyone originally sits at the bargaining table, with costs accidentally spilling over. Colonial cotton was not spillage. Slaves could not refuse; Indian weavers faced Company gunboats and tariff policy. Calling systematic plunder a negative externality mistakes coercion for oversight. \textbf{Understanding a mechanism neither excuses it nor permits measuring it with the wrong ruler.}

Second, externalities expose false prices without deciding correct ones. What is a tonne of smoke worth? A childhood? How many pounds compensate Sarah's afternoons without daring to sing? Economics estimates the first two but cannot calculate the third, not through imprecision but because the question lies outside arithmetic. The tool marks a ledger's boundary; what stands beyond requires other protections.

\section{Chimneys Move; Chains Remain}
If these seem two-century-old debts, examine your T-shirt again.

Cotton grows in Uzbekistan or the American South, yarn is spun in Vietnam, cloth woven in China and garments sewn in Dhaka before sale in Paris or your hometown. Global chains enlarge Manchester's model, placing the cheapest-billed stage farthest from consumers. In 2013, an eight-storey garment building outside Dhaka collapsed, killing over a thousand. Many had feared entering after cracks appeared the previous day, but wage deductions pushed them inside. You probably recognise brands made there. This is no isolated brand exception: externality changed time zones. Chimneys disappeared from view because they moved.

Time discipline changed skin too. Factory clocks and fine schedules governed Annie; phone countdowns and algorithmic orders govern couriers. A rider runs a red light not through ignorance of traffic rules, but because the workshop prices every minute. Systems calculate even thirty seconds waiting for a lift, more precisely than foremen. Luddites broke wage-devouring machines. Nobody breaks algorithms, because no single visible, tangible machine embodies them.

Industrialisation's largest externality bill is now before us. Most atmospheric carbon dioxide comes from over two centuries of coal and oil burning; early industrialisers contributed the largest historical emissions, while floods, drought and rising seas often strike the poorest, least-emitting countries first. An absent nineteenth-century ledger line dominates twenty-first-century negotiations. A local growth-and-pollution problem becomes a planetary dispute over who pays, in what shares and against which security.

Some things genuinely changed and must be acknowledged. An ordinary Chinese teenager today enjoys clothing, food, medicine and information exceeding those of Europe's richest people two centuries ago. Compulsory schooling replaced child labour through those centimetre-by-centimetre advances. The question is not whether industrialisation brought good things, but \textbf{whether distributing them still requires someone else to pay invisible bills}.

\section{Your Turn: Who Should Pay?}
Return to the 19.90-yuan T-shirt.

Suppose someone makes an honest T-shirt: fair cotton-farmer income, eight-hour mill days, fully treated wastewater and carbon costs included. All externalities return to the books. It costs ninety-nine yuan.

Answer with reasons: \textbf{who should pay the additional eighty yuan?}

Consumers? The case is strong: beneficiaries should pay; buying two fewer helps workers and the planet. But the cheapest buyers often have the fewest options: a child counting pocket money, a family relying on inexpensive winter clothing. Why send them the moral invoice first?

Brands and companies? They took profits and generated externalities. Yet they reply that competition thins margins to paper; a one-yuan increase sends customers next door. Are consumers not still deciding?

The state, forcing industry-wide higher prices through law? The 1833 Factory Act did something similar despite false predictions of collapse. But states compete. Raise standards and factories move elsewhere: that is how Dhaka's building arose.

Hear a fourth voice too: perhaps the question is wrong. Instead of who pays eighty yuan, ask why clothing must be designed for one season. Cheapness sometimes means deliberate physical and aesthetic short lives, encouraging replacement. This shifts attention from who pays the bill to who designs it.

When Sarah Gooder dared not sing underground, nobody thought this bill needed paying. You can see it now, and seeing begins every argument. What will you do with the next T-shirt? The answer is absent from textbooks and this chapter. It lies in your basket and your reasons.
